% answers_short.tex —— 2003 年·全国新课程卷·文科 · 简短答案速查
% 由 tools/make_short_answers.py 自 answers.tex 的 \gksolhead/\gkanswer 提取，
% 并与 2003/2003_数学真题_全国新课程文_答案要点.md 逐题比对后生成，勿手改（改源 MD 后重跑）。

\begin{gknotice}
\textbf{使用说明}：本节只给最终答案便于快速核对，完整推导见解析版。答案为编写者独立推导，\textbf{非官方标准答案}。
\end{gknotice}

\gkpart{选择题}
\begin{center}\begin{tabular}{cccccccccccccc}
\toprule
题号 & 1 & 2 & 3 & 4 & 5 & 6 & 7 & 8 & 9 & 10 & 11 & 12 & 13 \\
\midrule
答案 & C & D & B & D & C & B & D & B & B & C & B & C & A \\
\bottomrule
\end{tabular}\end{center}

\gkpart{填空题}
\begin{itemize}[leftmargin=2.6em]
  \item[\textbf{14.}] $-\dfrac{21}{2}$
  \item[\textbf{15.}] $6,30,10$
  \item[\textbf{16.}] $S_{BCD}^2=S_{ABC}^2+S_{ACD}^2+S_{ADB}^2$
  \item[\textbf{17.}] $48$
\end{itemize}

\gkpart{解答题}
\begin{itemize}[leftmargin=2.6em]
  \item[\textbf{18.}] (1) $\dfrac{2}{\sqrt3}$；(2) 见解析版
  \item[\textbf{19.}] (1) $a=-\dfrac12$，公切线 $y=x-\dfrac14$；(2) 见解析版
  \item[\textbf{20.}] (1) $a_2=4,\ a_3=13$；(2) $a_n=\dfrac{3^n-1}{2}$
  \item[\textbf{21.}] (1) $0.17575$；(2) $0.012$
  \item[\textbf{22.}] $\varphi=\dfrac\pi2$，$\omega=\dfrac23$ 或 $2$。
  \item[\textbf{23.}] 见解析版；轨迹为椭圆，焦点分类按 $a$ 与 $1/\sqrt2$ 的大小确定。
\end{itemize}
